%% kapa: exam=1 round=27 year=2016 period=1 form=A subject=civil_law subject_ko=민법 seq=5 number=5
%% source: Q-Net 감정평가사 기출문제 (cst00309) artlSeq=5169402 file=extracted/1008/5169402/제27회 감정평가사 1차 문제/제27회 감정평가사 1교시문제(A형).pdf
%% conversion: kordoc 4.24.51 -> tools/gen1.py (auto); review=proofread-sonnet; flags=none
\begin{question}{5}
주물과 종물에 관한 설명으로 옳은 것은? (다툼이 있으면 판례에 따름)
\begin{choices}
  \item 독립한 부동산은 종물이 될 수 없다.
  \item 주물을 처분할 때 당사자의 특약으로 종물만을 별도로 처분할 수도 있다.
  \item 주물 위에 설정된 저당권의 효력은, 법률의 규정 또는 다른 약정이 없으면, 설정 후의 종물에까지 미치지 않는다.
  \item 구분건물의 전유부분에 대한 가압류 결정의 효력은 특별한 사정이 없는 한 그 대지권에 미치지 않는다.
  \item 권리 상호간에는 주물과 종물의 법리가 적용되지 않는다.
\end{choices}
\end{question}
\answer{2}
